\chapter{Conclusion}
\label{chap:conclusion}


\iffalse{
Your conclusion actually follows the same structure as your abstract.

\begin{itemize}
	\item The domain, i.e., explain the application domain in which your thesis is applicable.
	\item The problem, i.e., motivate what the problem is you are solving.
	\item The method, i.e., describe the method you have proposed, and, compare it, briefly \& high level to other approaches. In particular, focus on the benefits of your approach versus the other approach(es)
	\item The evaluation, i.e., how did you evaluate your method, what results did you obtain  and what do the results indicate?
\end{itemize}

However, you present what the reader has seen, again, using the same concepts as you have already used in your abstract.

\paragraph{Future Work.}
Furthermore, you present interesting directions for future work!

}\fi


This thesis explored the application of object-centric process mining to life cycle assessment.
It contributes to the automation of LCI analysis, a critical phase in LCA responsible for modeling and quantifying the inputs and outputs of a system.
As a particular use case, we focused on waste detection, as waste can be derived from the inventory.
We motivated our research by highlighting that traditional LCI analysis is often a manual, time-consuming, and costly procedure and typically results in static models that are difficult to update.




For this purpose, we proposed the object flow miner, a discovery algorithm that constructs an object flow net directly from object-centric event logs.
These nets visualize the transformation of input objects into output objects for every activity, together with the corresponding object types and multiplicities.
In addition, the inputs and outputs of the entire process are marked.
The advantage of this approach is its ability to distinguish between different configurations of the same activity.
In contrast to existing process mining algorithms that focus on control flow, the OFN explicitly visualizes the object flow based on inputs and outputs.
We further extended our approach by annotating an OFN with waste data derived from quantitative mass attributes or object counts, achieving waste detection.



As a prerequisite, we required high-quality OCELs with certain restrictions, such as a total order of events, and the absence of object attribute changes.
Optionally, filtering parameters and the desired waste detection can also be specified.
On these inputs, the role functions are computed heuristically, which are used for the object flow miner that discovers an OFN and annotates it with waste.
This effectively automates the LCI analysis.



    
The evaluation of our approach showed its effectiveness in both theoretical and practical scenarios.
The algorithm discovered the basic constructs such as exclusive choice, loops, and non-free choice found in process models, but the discovery of parallelism is poor, and the OFN semantics do not allow for discovering implicit places.
Our method generated as-is models of the process with close to perfect fitness on four public OCELs.
When applied, filtering successfully reduced the size of process models to concise representations of the underlying processes.
The application to the hinge production log validated the applicability of the method for waste detection.
The annotated OFN revealed waste information on object types, as well as implicit waste on activities.
This information is derived from a provided mass attribute or counting automatically.
By applying the waste detection in reverse to the input object types as well, our approach semi-automated the LCI analysis phase, requiring only two manual hints.
The runtime performance was shown to be acceptable for practical use.
However, a major limitation is that we did not compare our approach with a reference LCI analysis on real-world data.



\section{Future Work}

The limitations of our algorithm suggest various directions for future work.
Our approach assumes that events never occur simultaneously, which is typically not the case, as we observed in our evaluation.
Therefore, handling simultaneous events would allow using the method on more OCELs.
Additionally, supporting object attribute changes would also make the approach more general, as currently attributes are assumed to be static.
The heuristic discovery of role functions, which is currently a case distinction based only on the preceding and succeeding activity of the object traces, could be improved by analyzing a wider context of events.
Further research into detecting parallelism and implicit places would enhance the precision of the discovered object flow nets.
Finally, visualizing the paths of individual objects through the model would allow for a deeper root-cause analysis of the detected waste.

